\Needspace{10\baselineskip}
\begin{center}\emph{The residues $+2$ and $-2$.}\end{center}

\Article{112}
We move on to the residues $+2$ and $-2$.

{\rightskip=0pt plus 4em
Collecting from Table II all the primes for which $+2$ is a residue gives $7,17,23,31,41,47,71,73,79.89,97$. We readily notice that none of them have either the form $8n+3$ or the form $8n+5$.\par}

Let us see whether we can turn this observed pattern into a certainty.

First, every composite number of the form $8n+3$ or $8n+5$ must have a prime factor of one of the forms $8n+3$ or $8n+5$. For primes of the forms $8n+1$, $8n+7$ alone can produce only numbers of the form $8n+1$ or $8n+7$. So if our observation holds in general, there is no number at all of the form \SourceLeaf{90}$8n+3$, $8n+5$ for which $+2$ is a residue. Certainly no such number below $100$ has $+2$ as a residue. If there are such numbers beyond this bound, let the smallest be $=t$. Then $t$ has either the form $8n+3$ or $8n+5$; $+2$ is a residue of it, but a nonresidue of all smaller numbers of these forms. Put $2\equiv aa$ (mod.~$t$). We can always choose $a$ odd and $<t$ (for $a$ has at least two positive values smaller than $t$ whose sum is $=t$, so one is even and the other odd; see Articles 104 and 105). Then let $aa=2+tu$, or $tu=aa-2$. Now $aa$ has the form $8n+1$, so $tu$ has the form $8n-1$, and $u$ has the form $8n+3$ or $8n+5$, according as $t$ has the latter or the former form. But the equation $aa=2+tu$ also gives $2\equiv aa$ (mod.~$u$), \emph{that is,} $2$ is also a residue of $u$. We readily see that $u<t$, so $t$ is not the smallest counterexample to our observation, contrary to the hypothesis. Thus the pattern we found by inspection holds in general.

Combining this with the proposition of Article 111 gives the following theorems.

I. \emph{For every prime of the form $8n+3$, $+2$ is a nonresidue and $-2$ a residue.}

II. \emph{For every prime of the form $8n+5$, both $+2$ and $-2$ are nonresidues.}

\Article{113}
Similarly, inspection of Table II gives the following primes for which $-2$ is a residue: $3,11,17,19,41,43,59,67,73,83,89,97$\HistoricalNote{*)}{By considering $-2$ as the product of $+2$ and $-1$; see Article 111.}. None have the forms $8n+5$, $8n+7$; let us see whether this observation too can be established as a general theorem. As in the preceding article, every composite number of the form $8n+5$ or $8n+7$ has a prime factor of the form $8n+5$ or $8n+7$. So if our observation holds in general, $-2$ cannot be a residue of any number of the form $8n+5$ or $8n+7$. If such numbers exist, let the smallest be $=t$, and put $-2=aa-tu$. Choose $a$ odd and smaller than $t$, as above; then $u$ has the form $8n+5$ or $8n+7$, according as $t$ has the form $8n+7$ or $8n+5$. From $aa+2=tu$ and $a<t$ we readily find that \SourceLeaf{91}$u$ is also smaller than $t$. Finally, $-2$ is a residue of $u$ too; \emph{that is,} $t$ is not the smallest counterexample to our observation, contrary to the hypothesis. So $-2$ must be a nonresidue of every number of the forms $8n+5$, $8n+7$.

Combining this with the proposition of Article 111 gives these theorems:

I.\quad \emph{For every prime of the form $8n+5$, both $-2$ and $+2$ are nonresidues,} as we already found in the preceding article.

II.\quad \emph{For every prime of the form $8n+7$, $-2$ is a nonresidue and $+2$ a residue.}

In both proofs we could also have chosen an even value of $a$. We would then have had to distinguish the case where $a$ has the form $4n+2$ from the case where $a$ has the form $4n$. The argument proceeds just as above, without any difficulty.

\Article{114}
One case remains: the prime has the form $8n+1$. The preceding method does not handle this case, which requires quite different devices.

For the prime modulus $8n+1$, let $a$ be any primitive root. Then (Article 62) $a^{4n}\equiv-1\,(\operatorname{mod.}8n+1)$. This congruence can also be written $(a^{2n}+1)^2\equiv2a^{2n}\,(\operatorname{mod.}8n+1)$, or $(a^{2n}-1)^2\equiv-2a^{2n}$. So both $2a^{2n}$ and $-2a^{2n}$ are residues of $8n+1$. Since $a^{2n}$ is a square not divisible by the modulus, both $+2$ and $-2$ are residues as well (Article 98).

\Article{115}
It will be useful to add another proof of this theorem, related to the preceding one as the second proof of the theorem of Article 108 (Article 109) is related to the first (Article 108). Readers familiar with these matters will then see more easily that neither pair of proofs is as different as it may seem at first sight.

\Needspace{5\baselineskip}
I.\quad For any prime modulus of the form $4m+1$, among the smaller numbers $1,2,3\ldots4m$, there are $m$ that can be congruent to a fourth power, while the remaining $3m$ cannot.

This follows easily from the principles of the previous section, but a proof without them is not difficult either. We proved that, for such a modulus, $-1$ \SourceLeaf{92}is always a quadratic residue. So let $ff\equiv-1$. If $z$ is any number not divisible by the modulus, the fourth powers of the four numbers $+z,-z,+fz,-fz$ (which are clearly incongruent) are congruent to one another. The fourth power of a number congruent to none of these four cannot be congruent to their fourth powers: otherwise the congruence $x^4\equiv z^4$, of the fourth degree, would have more than $4$ roots, contrary to Article 43. Thus all the numbers $1,2,3,\ldots4m$ give only $m$ incongruent fourth powers; among those same numbers, $m$ are congruent to them, and the rest cannot be congruent to any fourth power.

\Needspace{5\baselineskip}
II.\quad For a prime modulus of the form $8n+1$, $-1$ can be congruent to a fourth power ($-1$ is a \emph{biquadratic residue} of that prime).

For the number of all biquadratic residues smaller than $8n+1$ (excluding zero) is $=2n$, \emph{that is,} even. Also, if $r$ is a biquadratic residue of $8n+1$, it is easy to prove that the value of the expression $\frac{1}{r}\,(\operatorname{mod.}8n+1)$ is one too. So all the biquadratic residues can be divided into classes as we divided the quadratic residues in Article 109; the rest of the proof proceeds in exactly the same way as there.

\Needspace{7\baselineskip}
III.\quad Now let $g^4\equiv-1$, and let $h$ be the value of the expression $\frac{1}{g}\,(\operatorname{mod.}8n+1)$. Then
\[
(g\pm h)^2=g^2+h^2\pm2gh\equiv g^2+h^2\pm2
\]
(because $gh\equiv1$). But $g^4\equiv-1$, so $-h^2\equiv g^4h^2\equiv g^2$, and finally $g^2+h^2\equiv0$ and $(g\pm h)^2\equiv\pm2$; \emph{that is,} both $+2$ and $-2$ are quadratic residues of $8n+1$. \mbox{\emph{Q. E. D.}}

\Needspace{9\baselineskip}
\Article{116}
The preceding results also readily give this general rule: $+2$ \emph{is a residue of every number divisible neither by $4$ nor by any prime of the form $8n+3$ or $8n+5$; it is a nonresidue of all the others} (\emph{for example,} all numbers of the forms $8n+3$, $8n+5$, whether prime or composite).

$-2$ \emph{is a residue of every number divisible neither by $4$ nor by any prime of the form $8n+5$ or $8n+7$, and a nonresidue of all the others.}

These elegant theorems were already known to the perceptive Fermat, \emph{Op. Mathem., p.} 168. But he never \SourceLeaf{93}communicated the proof he claimed to have. Later the distinguished Euler repeatedly sought it in vain; the distinguished Lagrange was the first to find a rigorous proof, \emph{Nouv. M\'em. de l'Ac. de Berlin}, 1775, \emph{pp.} 349, 351. Euler seems still to have been unaware of this when he wrote the paper preserved in \emph{Opusc. Analyt.}, Vol. I, p. 259.

\Needspace{10\baselineskip}
\begin{center}\emph{The residues $+3$ and $-3$.}\end{center}

\Article{117}
We move on to the residues $+3$ and $-3$, starting with the latter.

From Table II we find the following primes for which $-3$ is a residue: 3, 7. 13, 19, 31, 37, 43, 61, 67, 73, 79, 97. None have the form $6n+5$. We prove as follows that there are no primes of this form for which $-3$ is a residue beyond the table's range either. First, every composite number of the form $6n+5$ must have a prime factor of the same form. As long as no primes of the form $6n+5$ have $-3$ as a residue, no such composite numbers occur either. If there are such numbers beyond the table's range, let the smallest be $=t$, and put $-3=aa-tu$. Choosing $a$ even and smaller than $t$, we have $u<t$, and $-3$ is a residue of $u$. If $a$ has the form $6n\pm2$, then $tu$ has the form $6n+1$, so $u$ has the form $6n+5$. \mbox{\emph{Q. E. A.}} For we assumed that $t$ was the smallest counterexample to our observation. If $a$ has the form $6n$, then $tu$ has the form $36n+3$, so $\frac{1}{3}tu$ has the form $12n+1$, and $\frac{1}{3}u$ has the form $6n+5$. But $-3$ is clearly also a residue of $\frac{1}{3}u$, and $\frac{1}{3}u<t$, \mbox{\emph{Q. E. A.}} Thus $-3$ cannot be a residue of any number of the form $6n+5$.

Every number of the form $6n+5$ has either the form $12n+5$ or $12n+11$, the former being included in $4n+1$ and the latter in $4n+3$. So we obtain the following theorems:

I.\quad \emph{For every prime of the form $12n+5$, both $-3$ and $+3$ are nonresidues.}

II.\quad \emph{For every prime of the form $12n+11$, $-3$ is a nonresidue and $+3$ a residue.}

\SourceLeaf{94}\Article{118}
From Table II we find the following numbers for which $+3$ is a residue: 3, 11, 13, 23, 37, 47, 59, 61, 71, 73, 83, 97. None have the form $12n+5$ or $12n+7$. That no numbers at all of the forms $12n+5$, $12n+7$ have $+3$ as a residue can be proved in exactly the same way as in Articles 112, 113 and 117, so we omit this work. Comparing with Article 111 gives the theorems:

I.\quad \emph{For every prime of the form $12n+5$, both $+3$ and $-3$ are nonresidues} (as we already found in the preceding article).

II.\quad \emph{For every prime of the form $12n+7$, $+3$ is a nonresidue and $-3$ a residue.}

\Needspace{8\baselineskip}
\Article{119}
This method gives no result for numbers of the form $12n+1$, which therefore require special devices. Inspection readily suggests that $+3$ and $-3$ are residues of every prime of this form. We need only prove that $-3$ is a residue of such numbers, since then $+3$ must be one too (Article 111). We will prove more generally that $-3$ is a residue of every prime of the form $3n+1$.

Let $p$ be such a prime, and let $a$, for the modulus $p$, belong to the exponent $3$ (such numbers exist by Article 54, since $3$ divides $p-1$). Then $a^3\equiv1\,(\operatorname{mod.}p)$, \emph{that is,} $a^3-1$, or $(a^2+a+1)(a-1)$, is divisible by $p$. But $a$ cannot be $\equiv1\,(\operatorname{mod.}p)$, since $1$ belongs to the exponent $1$. So $a-1$ is not divisible by $p$, but $a^2+a+1$ is, and therefore so is $4aa+4a+4$; \emph{that is,} $(2a+1)^2\equiv-3\,(\operatorname{mod.}p)$, or $-3$ is a residue of $p$. \mbox{\emph{Q. E. D.}}

This proof, which is independent of the preceding ones, also covers the primes of the form $12n+7$ already dealt with in the preceding article.

We also note that this argument could be presented along the lines of the method used in Articles 109 and 115, but for brevity we do not pursue this.

\SourceLeaf{95}\Article{120}
The preceding results readily give the following theorems (see Articles 102, 103 and 105).

I. $-3$ \emph{is a residue of all numbers divisible neither by $8$, nor by $9$, nor by any prime of the form $6n+5$, and a nonresidue of all the others.}

II. $+3$ \emph{is a residue of all numbers divisible neither by $4$, nor by $9$, nor by any prime of the form $12n+5$ or $12n+7$, and a nonresidue of all the others.}

Keep this particular case especially in mind:

$-3$ is a \emph{residue} of all primes of the form $3n+1$, or equivalently \emph{all those that are residues of $3$}; it is a \emph{nonresidue} of all primes of the form $6n+5$, or, excluding $2$, all those of the form $3n+2$, \emph{that is, all those that are nonresidues of $3$.} All the other cases clearly follow from this.

The propositions concerning the residues $+3$ and $-3$ were already known to Fermat, \emph{Opera Wallisii}, Vol. II, \emph{p.} 857. But the distinguished Euler first gave proofs, \emph{Comm. nov. Petr.}, Vol. VIII, \emph{p.} 105 ff. It is all the more surprising that proofs of the propositions concerning the residues $+2$ and $-2$, based on entirely similar devices, always escaped his insight. See also the paper by the distinguished Lagrange, \emph{Nouv. M\'em. de l'Ac. de Berlin}, 1775, \emph{p.} 352.
